\section{Positive-characteristic orbit density}\label{sec:orbits}

We prove that on each positive-characteristic height stratum, a stabilizer
orbit satisfies no nonzero analytic equation near its base point.

\begin{theorem}[Local orbit density]\label{orb:density}
Fix $1\leq h<n$ and a $C$-valued point $x$ of the parameter polydisc
$I_h=0$ with $u_h(x)\neq0$.
For $g\in A_n^\times$, write $y_g$ for its orbit point, with the
transport convention of Lemma~\ref{def:stabilizer}.
Let $G$ be any open subgroup of $A_n^\times$.  If an analytic germ $P$
at $x$ on the parameter polydisc defined by $I_h=0$ vanishes at all
points of the orbit $\{y_g:g\in G\}$ sufficiently close to $x$, then
$P=0$.

It suffices to assume that, for every $d\in A_n$,
$P(y_{1-p^Nd})=0$ for all sufficiently large $N$.
\end{theorem}

The proof will compose the Tate coordinates of
Section~\ref{sec:deformations} with comparison series whose coefficients
are analytic germs in the parameters.  These coefficients need not share
one domain of convergence.  Instead, at individual orbit points we will
choose convergent series with integral coefficients that match each
fixed finite initial segment eventually.  The following lemma shows
that this matching suffices to turn formal Taylor identities into
numerical estimates.

\begin{lemma}[Finite analytic jets]\label{orb:finite-jets}
Let $F$ be a complete nonarchimedean valued field, let $r\geq1$, and use
the maximum norm on $F^r$.  Let $\theta=(\theta_i)_{i=1}^r$ be an analytic germ at
$0\in F^r$ with $\theta(0)=0$, and let $a_j$, for $j\geq1$, be scalar
analytic germs at $0$.  Define formal series
\[
 U_i(X)=\sum_{j\geq1}a_j(X)\theta_i(X)^j
 \quad\text{in }F[[X_1,\ldots,X_r]],
\]
using the Taylor expansions of the germs.  Suppose that the linear part
$L=dU_0$ is invertible.

Let $\delta_N\to0$ in $F^r$, and suppose that coefficients
$a_{N,j}\in\OO_F$ have been chosen for every $N$ and $j\geq1$ so that,
for every $M\geq1$, all sufficiently large $N$ satisfy
\[
 a_{N,j}=a_j(\delta_N)\qquad(1\leq j\leq M).
\]
This includes the requirement that the finitely many germs in this
formula are defined at $\delta_N$.  For sufficiently large $N$, put
\[
 U_{N,i}=\sum_{j\geq1}a_{N,j}\theta_i(\delta_N)^j.
\]
These numerical series converge.  If $P$ is an analytic germ with
$P(\delta_N)=0$ eventually, write its unique formal expansion as
\[
 P(X)=\sum_{\ell\in\mathbb N^r}B_\ell U(X)^\ell.
\]
Then there is a constant $D>0$ such that
$\|\delta_N\|\leq D\|U_N\|$ for all sufficiently large $N$.
Moreover, for every $m\geq0$ there is $D_m>0$ such that
\begin{equation}\label{orb:finite-jet-estimate}
 \left|\sum_{|\ell|\leq m}B_\ell U_N^\ell\right|
 \leq D_m\|U_N\|^{m+1}
\end{equation}
for all sufficiently large $N$.  The threshold and $D_m$ may depend on
$m$.  No convergence of the full formal series $U$ or
$\sum_\ell B_\ell Z^\ell$ is required.
\end{lemma}

\begin{proof}
The formal definition of $U$ is valid because
$\theta_i\in(X_1,\ldots,X_r)$, so only finitely many summands contribute
to each total degree.  Formal inversion applies because $L$ is
invertible, and gives the asserted unique expansion of $P$.

For fixed $m\geq1$, the finite sum
\[
 F_{m,i}(X)=\sum_{j=1}^m a_j(X)\theta_i(X)^j
\]
is an actual analytic germ.  Its Taylor polynomial through total degree
$m$ is $J_mU_i$, the corresponding truncation of $U_i$.  Analyticity and
$\theta(0)=0$ give $\|\theta(\delta)\|\leq c\|\delta\|$ near $0$.
For sufficiently large $N$, coefficient matching and integrality give
\[
 \|U_N-F_m(\delta_N)\|
 \leq\|\theta(\delta_N)\|^{m+1}
 \leq c^{m+1}\|\delta_N\|^{m+1}.
\]
Here $\|\theta(\delta_N)\|<1$, which also proves convergence of the
numerical series defining $U_N$.  The Taylor estimate for the finitely
many analytic functions $F_{m,i}$ therefore yields
\begin{equation}\label{orb:jet-approximation}
 U_N=J_mU(\delta_N)+R_{m,N},\qquad
 \|R_{m,N}\|\leq c_m\|\delta_N\|^{m+1}.
\end{equation}
Explicitly, an analytic series $H=\sum_{|\beta|\geq m+1}h_\beta X^\beta$
converging on a closed polydisc of radius $\rho$ satisfies
$|H(\delta)|\leq M\rho^{-(m+1)}\|\delta\|^{m+1}$ for
$\|\delta\|\leq\rho$, where
$M=\sup_\beta |h_\beta|\rho^{|\beta|}<\infty$.

For $m=1$, write $U_N=L\delta_N+R_{1,N}$.  If
$\delta_N\neq0$ and $N$ is sufficiently large, then
\[
 \|R_{1,N}\|<\|L^{-1}\|^{-1}\|\delta_N\|
 \leq\|L\delta_N\|.
\]
The ultrametric inequality gives
$\|U_N\|=\|L\delta_N\|$, and hence
$\|\delta_N\|\leq\|L^{-1}\|\|U_N\|$.
For $\delta_N=0$ the same bound holds directly.  For sufficiently
large $N$, $U_N=0$ therefore implies $\delta_N=0$.

Put $Q_m(Z)=\sum_{|\ell|\leq m}B_\ell Z^\ell$ and $V_m=J_mU$.
The analytic germ $P-Q_m(V_m)$ has zero Taylor coefficients through
degree $m$: this follows from the formal identity for $P$ and
$V_m\equiv U\pmod{(X)^{m+1}}$.  Consequently,
\[
 |P(\delta_N)-Q_m(V_m(\delta_N))|
 \leq c'_m\|\delta_N\|^{m+1}.
\]
Both $U_N$ and $V_m(\delta_N)$ tend to zero.  On the unit polydisc the
finite polynomial $Q_m$ is Lipschitz: telescoping each monomial gives
the constant
$\max(1,\max_{1\leq|\ell|\leq m}|B_\ell|)$.
Equation~\eqref{orb:jet-approximation}, the preceding bound, and
$P(\delta_N)=0$ now imply
\[
 |Q_m(U_N)|\leq c''_m\|\delta_N\|^{m+1}.
\]
The inverse norm estimate proves
Equation~\eqref{orb:finite-jet-estimate} for $m\geq1$.
For $m=0$, continuity gives $B_0=P(0)=0$, so the estimate also holds.
Every use of the formal identity in this argument has taken place
after a finite truncation.
\end{proof}

\begin{proof}[Proof of Theorem~\ref{orb:density}]
Put $r=n-h$ and $q=p^h$.  Choose a $\Zp$-basis
$z_1,\ldots,z_r$ of $T_x$ as in
Lemma~\ref{def:tate-kernel}.  Proposition~\ref{def:coordinates} says
that
\[
 \theta(y)=\bigl(\Theta_y(z_i)\bigr)_{i=1}^r,
 \qquad \theta(x)=0,
\]
is analytic with invertible differential at $x$.  Fix the formal group
isomorphism
\[
 \alpha_x:H_x\xrightarrow{\sim}\Gamma_h,
 \qquad \alpha_x(T)=\sum_{j\geq1}a_j^xT^j,
\]
of Proposition~\ref{def:natural-comparison}.  Its forward coefficients
belong to $\OO_C$, its linear coefficient is nonzero, and
$\nu_x=\alpha_x\Theta_x$ is a nonzero natural map
$\mathcal H_n\to\mathcal H_h$.

\smallskip
\noindent\emph{Formal coordinates.}
Write $[p]_{H_y}(T)=f_y(T^q)$ and $b(y)=f'_y(0)$, so that
$b(x)\neq0$.  By Lemma~\ref{def:recursion}, a series
$\alpha(T)=\sum_{j\geq1}A_jT^j$ intertwines this multiplication by
$p$ with $q$th powers exactly when
\begin{equation}\label{orb:coefficient-germs}
 A_j^q-b(y)^jA_j
 =\sum_{i<j}A_i[Z^j]f_y(Z)^i
 \qquad(j\geq1).
\end{equation}
Inductively choose an analytic germ $a_j(y)$ satisfying this equation
and $a_j(x)=a_j^x$.  The derivative in the new variable is
$-b(x)^j\neq0$, so the analytic implicit function theorem applies.
One can see the analytic branch directly: with the earlier coefficients
fixed and $A_j=a_j^x+t$, the equation becomes
\[
 t=b(y)^{-j}\bigl(t^q+e(y)\bigr),\qquad e(x)=0.
\]
On a sufficiently small closed parameter polydisc, choose $\eta>0$
with $\|b^{-j}\|\eta^{q-1}<1$ and
$\|b^{-j}e\|\leq\eta$.  The right side is a strict contraction on
the ball of analytic functions of norm at most $\eta$; its fixed point
has value zero at $x$.  This proves the required germ at every finite
stage.

Set $\delta=y-x$, where this subtraction is ordinary parameter
subtraction.  Define
\[
 U_i(\delta)=\sum_{j\geq1}a_j(x+\delta)
                   \Theta_{x+\delta}(z_i)^j
 \quad\text{in }C[[\delta_1,\ldots,\delta_r]].
\]
Each coefficient of $\delta$ uses only finitely many $a_j$.
The linear matrix is $a_1^x d\theta_x$, which is invertible.
Thus, if $P$ were nonzero, there would be a fixed family
$(B_\ell)_{\ell\in\mathbb N^r}$ in $C$, not all zero, such that
\[
 P(x+\delta)=\sum_\ell B_\ell U(\delta)^\ell
 \quad\text{formally}.
\]
The germs $a_j$, the formal coordinates $U$, and the coefficients
$B_\ell$ are now fixed.  No common analytic domain for all $a_j$ has
been chosen.

\smallskip
\noindent\emph{Matching actual coefficients on an orbit.}
Fix $d\in A_n$.  For all sufficiently large $N$, the element
$g_N=1-p^Nd$ is a unit in $G$, and its orbit point $y_N=y_{g_N}$
tends to $x$ by the continuity discussed before Lemma~\ref{def:stabilizer}.  The subtraction in
$g_N$ is in the endomorphism ring.  Transport sends $g_Nz_i$ into
$T_{y_N}$; applying the group homomorphism $\Theta_{y_N}$ therefore
gives
\begin{equation}\label{orb:transport-powers}
 \Theta_{y_N}(z_i)
 =[p]_{H_{y_N}}^N\bigl(\Theta_{y_N}(dz_i)\bigr).
\end{equation}
This identity is obtained using formal-group addition.

We next choose a numerical series at each $y_N$.  For every $M\geq1$
choose an integer $N_M\geq M$, strictly increasing with $M$, so that
for every $N\geq N_M$ the following conditions hold: the first $M$
germs $a_j$ are defined at $y_N$, their instances of
Equation~\eqref{orb:coefficient-germs} hold there, and $b(y_N)$ and $a_1(y_N)$
are nonzero.  Such thresholds exist because $y_N\to x$.  For
$N\geq N_1$ put
\[
 K(N)=\max\{M\leq N:N_M\leq N\}.
\]
Then $K(N)\to\infty$.  Prescribe
$a_{N,j}=a_j(y_N)$ for $1\leq j\leq K(N)$, and extend the sequence
inductively by choosing roots in $C$ of
Equation~\eqref{orb:coefficient-germs} at $y_N$.  Lemma~\ref{def:recursion}
shows that every chosen coefficient, including the prescribed ones,
belongs to $\OO_C$.  The resulting actual series
\[
 \alpha_N(T)=\sum_{j\geq1}a_{N,j}T^j
\]
converges on $\mm_C$ and satisfies the full identity
\[
 \alpha_N\bigl([p]_{H_{y_N}}(T)\bigr)=\alpha_N(T)^q.
\]
This intertwining identity is the only property of $\alpha_N$ needed
below.  In particular, for every fixed $j$ its coefficient satisfies
$a_{N,j}=a_j(y_N)$ eventually and hence $a_{N,j}\to a_j^x$.

Define convergent numerical quantities
\[
 U_{N,i}=\alpha_N\bigl(\Theta_{y_N}(z_i)\bigr),\qquad
 s_{N,i}=\alpha_N\bigl(\Theta_{y_N}(dz_i)\bigr).
\]
Iterating the intertwining identity in
Equation~\eqref{orb:transport-powers} gives the exact coordinate
identity
\begin{equation}\label{orb:rooted-coordinates}
 U_N=s_N^{q^N}.
\end{equation}
All evaluations are on topologically nilpotent elements; the integral
series involved preserve this domain.

Continuity in the parameter gives
$\Theta_{y_N}(dz_i)\to\Theta_x(dz_i)$.  Choose $c<1$ bounding the
norms of these finitely many arguments for all sufficiently large $N$
and of their limits.  Integrality bounds the tails of $\alpha_N$ past
degree $J$ by $c^{J+1}$, uniformly in $N$; the same bound holds for
$\alpha_x$.  The first $J$ coefficients converge term by term.  Letting
first $N\to\infty$ and then $J\to\infty$ proves
\begin{equation}\label{orb:limit-tuple}
 s_N\longrightarrow s(d):=\bigl(\nu_x(dz_i)\bigr)_{i=1}^r
 \quad\text{in }\mm_C^r.
\end{equation}
This argument also covers a zero coordinate, or a zero limiting tuple.

\smallskip
\noindent\emph{Frobenius roots and one formal equation.}
Apply Lemma~\ref{orb:finite-jets} with
$\delta_N=y_N-x$, the germs $a_j(x+\delta)$, and the actual
coefficients $a_{N,j}$.  The vanishing hypothesis gives
$P(y_N)=0$ eventually.  Thus, for every fixed $m\geq0$,
\[
 \left|\sum_{|\ell|\leq m}B_\ell U_N^\ell\right|
 \leq D_m\|U_N\|^{m+1}
\]
eventually.  The constant and threshold may depend on the fixed $d$
as well as on $m$.  Since $C$ is perfect of characteristic $p$,
Equation~\eqref{orb:rooted-coordinates} permits taking the unique
$q^N$th root of this finite sum, giving
\[
 \left|\sum_{|\ell|\leq m}B_\ell^{1/q^N}s_N^\ell\right|
 \leq D_m^{1/q^N}\|s_N\|^{m+1}.
\]
For this fixed $m$, the finitely many numbers
$B_\ell^{1/q^N}$ are uniformly bounded.  By
Equation~\eqref{orb:limit-tuple}, replacing $s_N$ by $s(d)$ in the
finite sum changes its value by a quantity tending to zero.
As $D_m^{1/q^N}\to1$, the ultrametric inequality gives the
full-sequence estimate
\begin{equation}\label{orb:rooted-limit}
 \limsup_{N\to\infty}
 \left|\sum_{|\ell|\leq m}B_\ell^{1/q^N}s(d)^\ell\right|
 \leq\|s(d)\|^{m+1}.
\end{equation}
When $s(d)=0$, the right side of the preceding numerical bound and
the replacement error both tend to zero, so the same conclusion
holds.

For each $d\in A_n$, Equation~\eqref{orb:rooted-limit} holds for every
$m\geq0$.  Lemma~\ref{formal:separation} therefore gives
\[
 \sum_\ell\lambda(B_\ell)s(d)^\ell=0
\]
for every $k$-linear functional $\lambda:C\to k$ and every $d\in A_n$.
Choose one such functional nonzero on a fixed nonzero coefficient
$B_{\ell_0}$.  Since the family $B_\ell$ was fixed before $d$, this
produces a single nonzero formal series over $k$ vanishing on the whole
set
\[
 \left\{\bigl(\nu_x(dz_i)\bigr)_{i=1}^r:d\in A_n\right\}.
\]
The $z_i$ are $\Qp$-linearly independent, $r=n-h<n$, and $\nu_x$ is
a nonzero natural map.  Proposition~\ref{bundles:nondegeneracy}
excludes precisely this formal equation.  The contradiction proves
the theorem.
\end{proof}
